\documentclass[11pt,a4paper]{amsart}
\pdfinfoomitdate=1
\pdftrailerid{}
\pdfsuppressptexinfo=15
\usepackage[T1]{fontenc}
\usepackage{lmodern}
\usepackage[margin=27mm]{geometry}
\usepackage[expansion=false]{microtype}
\usepackage{amsmath,amssymb,mathtools}
\usepackage{array,booktabs,float}
\usepackage{xcolor}
\usepackage[numbers,sort&compress]{natbib}
\usepackage[colorlinks=true,linkcolor=blue!55!black,citecolor=blue!55!black,urlcolor=blue!55!black]{hyperref}
\usepackage[nameinlink,noabbrev]{cleveref}
\raggedbottom

\newtheorem{theorem}{Theorem}[section]
\newtheorem{proposition}[theorem]{Proposition}
\newtheorem{corollary}[theorem]{Corollary}
\newtheorem{lemma}[theorem]{Lemma}
\theoremstyle{definition}
\newtheorem{definition}[theorem]{Definition}
\theoremstyle{remark}
\newtheorem{remark}[theorem]{Remark}

\newcommand{\F}{\mathcal F}
\newcommand{\B}{\mathcal B}
\newcommand{\depth}{\operatorname{depth}}
\newcommand{\Fix}{\operatorname{Fix}}
\newcommand{\pern}{\operatorname{per}}

\title[Odd-run reversal]{Odd-Run Reversal on Cyclic Binary Words:\
Exact Recurrence and Sharp Parity-Dependent Transients}
\author{Anonymous}
\date{}
\subjclass[2020]{37P35, 68Q80, 05A15}
\keywords{finite dynamics, cyclic binary word, run, shrinking system,
preperiod, periodic-point zeta function}

\begin{document}

\begin{abstract}
On labelled cyclic binary words of length \(n\), consider the parallel rule
that flips every bit belonging to a maximal run of odd length.  A boundary
survives one round exactly when its two incident run lengths have the same
parity; no new boundary is created.  We prove that a word is recurrent
exactly when all its cyclic run lengths have one parity.  All eventual
periods are therefore at most two.  For odd \(n\), the two constant words
form the only recurrent orbit.  For \(n=2m\), the fixed-state count is
\(2^{m+1}-2\), while the number of period-two states is
\[
 p_{2m}=\sum_{\substack{2\le r\le2m\\r\ {\rm even}}}
 \frac{4m}{r}\binom{m+r/2-1}{r-1}.
\]
The main temporal result is the sharp maximum preperiod
\[
 \max\depth=
 \begin{cases}
 (n-1)/2,&n\ \text{odd},\\
 \lfloor(n-2)/4\rfloor,&n\ \text{even}.
 \end{cases}
\]
The odd case follows from run coalescence.  The even case needs a different
coordinate: label surviving boundaries by their site parity, delete a label
when its two neighbours differ, and use a realization cost that drops by at
least four at every mixed step.  Binary-run enumeration, shrinking-cell
models, cyclic compositions, and zeta bookkeeping receive zero contribution
credit.  The residual scope is only this rule-specific conjunction.
Ownership, priority, and external circulation remain on hold.
\end{abstract}

\maketitle

\section{Introduction}

A local update stated in terms of maximal runs is not a fixed-radius cellular
automaton: a decision can depend on the distance to the next change of bit.
It nevertheless has a simple lower-dimensional trace.  The rule studied here
flips every odd run simultaneously.  Instead of tracking the bits, we track
the boundaries between runs.  They can disappear but never appear.

This places the mechanism near shrinking cellular automata, where deleted
cells make their surviving neighbours adjacent
\cite{RosenfeldWuDubitzki1983,ModaneseWorsch2016,KutribMalcherWendlandt2017}.
Those models and their language-recognition theory are background, not claims
of this paper.  Static enumeration of binary runs and its connection with
compositions are also mature; a current systematic account is
\cite{BaladoSilvestre2026}.  We give all such run and composition machinery
zero contribution credit.

After that subtraction, four exact statements remain for one specified
finite self-map.  First, boundary survival is governed only by the parities
of the two incident run lengths.  Second, this gives a complete recurrent
classification and labelled census.  Third, elementary coalescence proves a
sharp clock at odd circumference.  Fourth, even circumference carries an
intrinsic parity label on each boundary; its induced eroder and a cost-drop
lemma prove a different sharp clock.

The parity distinction is essential.  Counting boundaries gives the exact
odd bound but loses a factor of two in the even case.  Conversely, boundary
site parity is globally well defined only when the circumference is even.
The two proof routes therefore meet at the survival lemma and then separate.

We count labelled words, not rotation classes.  We do not determine every
basin layer or claim that the update is identical to a named shrinking
automaton.  A bounded owner search did not expose the same temporal
conjunction, but a search miss is neither novelty evidence nor a priority
certificate.  External dissemination remains \textbf{HOLD}.

\section{The update and its boundary trace}

Let \(\mathbb Z_n=\mathbb Z/n\mathbb Z\), with \(n\ge1\), and let
\(\mathcal W_n=\{0,1\}^{\mathbb Z_n}\).  A \emph{cyclic run} is a maximal
cyclic interval on which a word is constant.  A constant word has one run,
of length \(n\).

\begin{definition}[odd-run reversal]\label{def:update}
For \(w\in\mathcal W_n\), the word \(\F_n(w)\) is obtained by flipping every
bit in every odd-length cyclic run of \(w\), simultaneously, and retaining
every bit in an even-length run.
\end{definition}

For a nonconstant word, let \(\B(w)\subseteq\mathbb Z_n\) be its set of run
starts: \(i\in\B(w)\) when \(w(i-1)\ne w(i)\).  In cyclic order write the
run lengths as \(\ell_0,\ldots,\ell_{r-1}\).  Here \(r=|\B(w)|\) is even.

\begin{lemma}[boundary survival]\label{lem:survival}
No new run boundary is created.  The boundary between consecutive runs of
lengths \(\ell_{i-1}\) and \(\ell_i\) survives one update if and only if
\[
                    \ell_{i-1}\equiv\ell_i\pmod2.
\]
\end{lemma}

\begin{proof}
Every bit of one old run receives the same flip decision, so no boundary can
appear inside it.  Let the bit on the left of an old boundary be \(a\), and
the bit on the right be \(1-a\).  Put
\(\epsilon=\ell_{i-1}\bmod2\) and
\(\eta=\ell_i\bmod2\).  The two new endpoint bits are
\[
                  a+\epsilon,\qquad 1-a+\eta
                  \quad\text{in }\mathbb F_2.
\]
They remain different exactly when \(\epsilon=\eta\).
\end{proof}

\begin{theorem}[recurrent classification]\label{thm:recurrence}
A word is recurrent under \(\F_n\) if and only if all of its cyclic run
lengths have the same parity.  In that case it is fixed when all run lengths
are even and has exact period two when all run lengths are odd.  Every orbit
has eventual period at most two.
\end{theorem}

\begin{proof}
By \cref{lem:survival}, boundary sets decrease under inclusion.  If an orbit
returns to a previous word, its boundary set cannot have decreased anywhere
along the intervening segment.  Every boundary therefore survives the first
round of that segment.  Equality of the two incident parities at every
boundary propagates around the cyclic run list, so all run lengths have one
parity.

Conversely, if all runs are even, no bit flips and the word is fixed.  If all
runs are odd, every bit flips.  The resulting complement has the same run
decomposition, so the next update returns to the original word.  A binary
word never equals its bitwise complement, hence this period is exactly two.
Every finite orbit eventually becomes recurrent, which proves the last
claim.
\end{proof}

\section{Exact recurrent census}

Let \(f_n\) and \(p_n\) denote the numbers of fixed and exact-period-two
states, respectively.

\begin{theorem}[labelled recurrent counts]\label{thm:census}
For odd \(n\),
\[
                         f_n=0,\qquad p_n=2.
\]
For \(n=2m\),
\begin{align}
 f_{2m}&=2^{m+1}-2,                                      \label{eq:fixed}\\
 p_{2m}&=\sum_{\substack{2\le r\le2m\\r\ {\rm even}}}
 \frac{4m}{r}\binom{m+r/2-1}{r-1}.                       \label{eq:two}
\end{align}
\end{theorem}

\begin{proof}
A nonconstant binary cyclic word has an even number of runs.  If \(n\) is
odd, neither a sum of even run lengths nor a sum of an even number of odd
run lengths can equal \(n\).  Thus only the two constant words recur.  Their
single run has odd length, so they form a two-cycle.

Now let \(n=2m\).  For a nonconstant fixed word, all boundary gaps are even.
Equivalently, its nonempty boundary set is an even-cardinality subset of one
of the two parity classes in \(\mathbb Z_{2m}\).  Each parity class has
\(m\) sites and \(2^{m-1}-1\) nonempty even subsets.  Each boundary set
supports two complementary words.  Adding the two constant words gives
\[
       2\cdot2(2^{m-1}-1)+2=2^{m+1}-2.
\]

For an exact-period-two word with \(r\) runs, every run length is positive
and odd.  Writing \(\ell_i=2x_i+1\), stars and bars gives
\[
 \#\{(\ell_0,\ldots,\ell_{r-1}):\ell_i\text{ odd},
       \ \sum_i\ell_i=2m\}
 =\binom{m+r/2-1}{r-1}.
\]
Choose a labelled distinguished boundary in \(2m\) ways and the bit after it
in two ways, then divide by the \(r\) possible distinguished boundaries of
the same word.  Summing over even \(r\) proves \eqref{eq:two}.
\end{proof}

For a finite map \(F\), write
\[
 \zeta_F(z)=\exp\!\left(\sum_{j\ge1}
          \frac{|\Fix(F^j)|}{j}z^j\right).
\]
This periodic-point construction is standard \cite{ArtinMazur1965}.

\begin{corollary}[finite zeta function]\label{cor:zeta}
\[
 \zeta_{\F_n}(z)=(1-z)^{-f_n}(1-z^2)^{-p_n/2},
\]
with \(f_n,p_n\) from \cref{thm:census}.
\end{corollary}

\begin{proof}
\Cref{thm:recurrence} leaves only fixed orbits and two-cycles.  The formula
is the routine product of one factor for every primitive cycle.
\end{proof}

\begin{table}[H]
\centering
\caption{Exact recurrent census and sharp maximum preperiod for small
orders.  The table is illustrative; the formulas hold for every \(n\).}
\label{tab:census}
\begin{tabular}{rrrrr}
\toprule
\(n\)& fixed states & period-two states & two-cycles & max preperiod\\
\midrule
1&0&2&1&0\\
2&2&2&1&0\\
3&0&2&1&1\\
4&6&10&5&0\\
5&0&2&1&2\\
6&14&32&16&1\\
7&0&2&1&3\\
8&30&90&45&1\\
9&0&2&1&4\\
10&62&242&121&2\\
\bottomrule
\end{tabular}
\end{table}


\section{The sharp odd clock}

Define the preperiod
\[
 \depth(w)=\min\{t\ge0:\F_n^t(w)\text{ is recurrent}\}.
\]

\begin{theorem}[odd circumference]\label{thm:odd-depth}
If \(n\) is odd, then
\[
                   \max_{w\in\mathcal W_n}\depth(w)=\frac{n-1}{2}.
\]
\end{theorem}

\begin{proof}
At a nonrecurrent state, some adjacent run parities differ.  The number of
deleted cyclic boundaries is even and positive, so at least two boundaries
disappear.  An odd cyclic word has at most \(n-1\) boundaries: a boundary at
every site would be an alternating cycle, which exists only at even
circumference.  Monotone boundary loss therefore gives
\(\depth(w)\le(n-1)/2\).

For sharpness, write \(n=2t+1\).  Use the cyclic run composition consisting
of one part of length two and \(2t-1\) parts of length one.  The unique even
run has an odd run on each side.  At the next update both incident
boundaries disappear, and these three runs coalesce into one even run.
Every other boundary, lying between two odd runs, survives.  The same
description repeats for the first \(t-1\) rounds, with two fewer odd
singleton runs each time.  The remaining run composition is then
\((2t,1)\); its last two boundaries disappear together, leaving the odd
constant run of length \(2t+1\).  Thus the first recurrent state occurs
after exactly \(t\) rounds.  (For \(t=1\), this last two-run collapse is the
only round.)  For \(n=1\), the constant word is already recurrent and the
same formula gives zero.
\end{proof}

\section{The parity eroder and the sharp even clock}

Assume throughout this section that \(n\) is even.  If a nonconstant word
has surviving boundaries
\[
                  b_0,b_1,\ldots,b_{r-1}\in\mathbb Z_n
\]
in cyclic order, define its \emph{boundary-parity word}
\[
                         q_i=b_i\bmod2.
\]
This label is intrinsic because reduction modulo two is well defined on
\(\mathbb Z_n\) exactly when \(n\) is even.  The run between \(b_i\) and
\(b_{i+1}\) has parity \(q_i+q_{i+1}\) in \(\mathbb F_2\).  Applying
\cref{lem:survival} at \(b_i\) yields the following autonomous rule.

\begin{definition}[boundary-parity eroder]\label{def:eroder}
For a cyclic binary word \(q\), let \(Dq\) be the subsequence that retains
position \(i\) precisely when
\[
                            q_{i-1}=q_{i+1}.
\]
The cyclic order of retained positions is unchanged.
\end{definition}

Thus the boundary-parity word after one \(\F_n\)-update is exactly \(Dq\).
The recurrent parity words are constant and alternating: they correspond,
respectively, to all-even and all-odd original run lengths.

For a nonempty cyclic word \(q\), let
\[
 e(q)=|\{i:q_i=q_{i+1}\}|,\qquad C(q)=|q|+e(q).
\]
Put \(C(\varnothing)=0\).

\begin{lemma}[realization cost]\label{lem:cost}
If \(q\) is the boundary-parity word of a cyclic word of circumference
\(n\), then \(q\) has even length and \(n\ge C(q)\).  Conversely, every
nonempty cyclic binary word \(q\) of even length is realizable as a
boundary-parity word at circumference \(C(q)\), and at every circumference
\(C(q)+2j\), \(j\ge0\).  The empty word is realized by either constant word.
\end{lemma}

\begin{proof}
The positive gap from \(b_i\) to \(b_{i+1}\) is even when
\(q_i=q_{i+1}\), so it costs at least two, and is odd otherwise, so it costs
at least one.  Summing the gap minima gives
\[
                 2e(q)+(|q|-e(q))=C(q).
\]
Taking every gap at its minimum realizes equality.  Adding two to any one
gap preserves every boundary parity and realizes all larger circumferences
of the same parity.  Since a binary cyclic word has an even number of
boundaries, alternating bit values can be assigned consistently between
the chosen positions.
\end{proof}

\begin{lemma}[four-unit cost drop]\label{lem:drop}
If \(q\) is neither constant nor alternating and has even length, then
\[
                            C(Dq)\le C(q)-4.
\]
\end{lemma}

\begin{proof}
Decompose \(q\) into cyclic constant runs.  A singleton run survives.  A
run of length two disappears.  A run of length at least three loses its two
endpoints and retains its interior.

Let \(a\) be the number of nonsingleton runs and \(h\le a\) the number of
length-two runs.  If \(s\) is the number of transitions of \(q\), then
\(e(q)=|q|-s\), so \(C(q)=2|q|-s\).  The update deletes exactly \(2a\)
symbols.  Only a disappearing length-two run can merge its two neighbouring
surviving runs; deleting one such run lowers the transition count by at most
two.  Thus, with \(s'\) the transition count after deletion (and with the
transition count of the empty word defined to be zero),
\[
 |Dq|=|q|-2a,\qquad s'\ge s-2h.
\]
Consequently
\[
 C(q)-C(Dq)=4a+s'-s\ge4a-2h.
\]
If \(a\ge2\), the right side is at least \(2a\ge4\).

It remains to consider \(a=1\).  All other constant runs are singletons.
Both the total length and the number of cyclic runs are even, so the unique
nonsingleton run has odd length, necessarily at least three.  It remains
nonempty, no run disappears, and \(s'=s\).  The cost drop is then exactly
four.
\end{proof}

\begin{theorem}[even circumference]\label{thm:even-depth}
If \(n\) is even, then
\[
        \max_{w\in\mathcal W_n}\depth(w)
        =\left\lfloor\frac{n-2}{4}\right\rfloor.
\]
\end{theorem}

\begin{proof}
A mixed even-length parity word has at least four symbols.  Its number of
equal adjacencies is positive and even, hence at least two.  Therefore its
cost is at least six.  If an original word has preperiod \(t\ge1\), its
parity words before rounds \(0,\ldots,t-1\) are mixed.  Applying
\cref{lem:drop} for the first \(t-1\) transitions and then
\cref{lem:cost} gives
\[
                       n\ge6+4(t-1)=4t+2.
\]
Thus \(t\le\lfloor(n-2)/4\rfloor\).

For sharpness, fix \(t\ge1\) and take
\[
                              q=0^{\,2t+1}1.
\]
Its cost is \(4t+2\).  Each eroder round removes the two endpoints of the
long zero run.  After exactly \(t\) rounds it becomes the alternating word
\(01\), and no earlier state is recurrent.  By \cref{lem:cost}, this
boundary word is realized by an original cyclic word at circumference
\(4t+2\); adding two to one gap realizes circumference \(4t+4\) with the
same parity dynamics.  These are exactly the two even circumferences for
which \(\lfloor(n-2)/4\rfloor=t\).  At \(n=2,4\), the claimed maximum is
zero and the recurrent classification verifies it directly.
\end{proof}

\section{Exact controls, scope, and conclusion}

The accompanying deterministic verifier exhausts all \(2^n\) labelled words
for \(1\le n\le16\).  It directly checks one-step boundary survival and the
induced parity eroder, as well as recurrence, the exact census, both sharp
depth formulas, and extremal witnesses.  A second exact lane enumerates all
\(349{,}524\) even-length parity words through length eighteen, constructs
their minimum-cost realizations, and checks the four-unit drop on all
\(349{,}488\) mixed cases.  A fresh run executes \(1{,}529{,}158\) exact
assertions.  This finite calculation is falsification and regression
evidence; none of the all-\(n\) statements depends on it.

The result is deliberately narrow.  It gives no complete transient-layer
distribution and no orbit census modulo rotation.  Shrinking automata,
binary-run enumeration, cyclic compositions, and periodic-point zeta
bookkeeping are prior tools and receive zero credit.  The labelled census is
a routine enumerative corollary once the temporal classification is known.
What remains is the exact map-specific conjunction for odd-run reversal: the
boundary survival law, complete recurrence, and the two parity-sensitive
sharp clocks.  Owner clearance, priority, and external release remain on
hold.

\bibliographystyle{plainnat}
\bibliography{references}

\end{document}
